% !TeX TS-program = pdflatex


\documentclass[a4paper]{article}

% \usepackage[default]{fontsetup}

\usepackage{fancyhdr}
\usepackage{extramarks}
\usepackage{amsmath}
\usepackage{amsthm}
\usepackage{amsfonts}
\usepackage{tikz}
\usepackage[plain]{algorithm}
\usepackage{algpseudocode}
\usepackage{enumerate}
\usepackage{tikz}

\usetikzlibrary{automata,positioning}

%
% Basic Document Settings
%  

\topmargin=-0.2in
\evensidemargin=0in
\oddsidemargin=0in
\textwidth=6.5in
\textheight=9.5in
\headsep=0.25in

\linespread{1.1}

\pagestyle{fancy}
\lhead{\hmwkAuthorName}
\chead{\hmwkClass : \hmwkTitle}
\rhead{\firstxmark}
\lfoot{\lastxmark}
\cfoot{\thepage}

\renewcommand\headrulewidth{0.4pt}
\renewcommand\footrulewidth{0.4pt}

\setlength\parindent{0pt}

%
% Create Problem Sections
%

\newcommand{\enterProblemHeader}[1]{
    \nobreak\extramarks{}{Problem \arabic{#1} continued on next page\ldots}\nobreak{}
    \nobreak\extramarks{Problem \arabic{#1} (continued)}{Problem \arabic{#1} continued on next page\ldots}\nobreak{}
}

\newcommand{\exitProblemHeader}[1]{
    \nobreak\extramarks{Problem \arabic{#1} (continued)}{Problem \arabic{#1} continued on next page\ldots}\nobreak{}
    \stepcounter{#1}
    \nobreak\extramarks{Problem \arabic{#1}}{}\nobreak{}
}

\newcommand*\circled[1]{\tikz[baseline=(char.base)]{
		\node[shape=circle,draw,inner sep=2pt] (char) {#1};}}


\setcounter{secnumdepth}{0}
\newcounter{partCounter}
\newcounter{homeworkProblemCounter}
\setcounter{homeworkProblemCounter}{1}
\nobreak\extramarks{Problem \arabic{homeworkProblemCounter}}{}\nobreak{}

%
% Homework Problem Environment
%
% This environment takes an optional argument. When given, it will adjust the
% problem counter. This is useful for when the problems given for your
% assignment aren't sequential. See the last 3 problems of this template for an
% example.
%

\newenvironment{homeworkProblem}[1][-1]{
    \ifnum#1>0
        \setcounter{homeworkProblemCounter}{#1}
    \fi
    \section{Problem \arabic{homeworkProblemCounter}}
    \setcounter{partCounter}{1}
    \enterProblemHeader{homeworkProblemCounter}
}{
    \exitProblemHeader{homeworkProblemCounter}
}

%
% Homework Details
%   - Title
%   - Class
%   - Due date
%   - Name
%   - Student ID

\newcommand{\hmwkTitle}{Homework\ \#03}
\newcommand{\hmwkClass}{Probability \& Statistics for EECS}
\newcommand{\hmwkDueDate}{To be announced}
\newcommand{\hmwkAuthorName}{}
\newcommand{\hmwkAuthorID}{}


%
% Title Page
%

\title{
    \vspace{2in}
    \textmd{\textbf{\hmwkClass:\\  \hmwkTitle}}\\
    \normalsize\vspace{0.1in}\small{Due\ on\ \hmwkDueDate\ at 23:59}\\
	\vspace{4in}
}

\author{
	Name: \textbf{\hmwkAuthorName} \\
	Student ID: \hmwkAuthorID}
\date{}

\renewcommand{\part}[1]{\textbf{\large Part \Alph{partCounter}}\stepcounter{partCounter}\\}

%
% Various Helper Commands
%

% Useful for algorithms
\newcommand{\alg}[1]{\textsc{\bfseries \footnotesize #1}}
% For derivatives
\newcommand{\deriv}[1]{\frac{\mathrm{d}}{\mathrm{d}x} (#1)}
% For partial derivatives
\newcommand{\pderiv}[2]{\frac{\partial}{\partial #1} (#2)}
% Integral dx
\newcommand{\dx}{\mathrm{d}x}
% Alias for the Solution section header
\newcommand{\solution}{\textbf{\large Solution}}
% Probability commands: Expectation, Variance, Covariance, Bias
\newcommand{\E}{\mathrm{E}}
\newcommand{\Var}{\mathrm{Var}}
\newcommand{\Cov}{\mathrm{Cov}}
\newcommand{\Bias}{\mathrm{Bias}}

\newcommand{\Prob}{P}

\begin{document}


% \maketitle
% \thispagestyle{empty}
% \pagebreak

\date{
Due date: Due on Oct. 11, 2026, 23:59 UTC+8}
\title{SI 140A.02  Probability \& Statistics for EECS, Fall 2026 \\
Homework 3}
\maketitle
Read all the instructions below carefully before you start working on the assignment, and before you make a submission.
\begin{itemize}
    \item You are required to write down all the major steps towards making your conclusions; otherwise you may obtain limited points of the problem.
    \item Write your homework in English; otherwise you will get no points of this homework.
    \item Any form of plagiarism will lead to $0$ point of this homework. 
    \item You are required to submit your homework in PDF format generated by LaTeX, otherwise you will get no points of this homework.
\end{itemize}
\newpage

\begin{homeworkProblem}[1]
% Source: Fall 2025 HW3, Problem 2(b), PDF page 3.
There are two states, Red and Blue, initially with 100 residents each. Every resident supports exactly one of two parties, Democrat or Republican. Initially, 60 residents of Blue and 30 residents of Red are Democrats. Ten residents move from Blue to Red, and nobody changes parties. Let $k$ be the number of Democrats among the ten people who move.

Whenever a resident is selected, the selection is uniform over all 200 people. Let $D$ be the event that the selected resident is a Democrat, and let $B$ be the event that the resident lives in Blue at the time of selection.
\begin{enumerate}[(a)]
\item Find $\Prob(D\mid B)$, $\Prob(D\mid B^c)$, and $\Prob(D)$ before the move. Does the move change $\Prob(D)$? Explain.
\item Determine all integer values of $k$ for which the proportion of Democrats increases strictly in \emph{both} states. Explain, using the law of total probability with the appropriate weights, how both proportions can increase while the overall proportion stays unchanged.
\item For this part, take $k=5$ and consider the population after the move. A survey interviews the selected resident only if their name appears on a list containing all Blue residents and all Democrats in Red. Let $A$ be the event that the selected resident is on this list. Find $\Prob(D\mid A)$ using the law of total probability conditional on $A$, with the partition $\{B,B^c\}$, and state the two group weights. Explain why this probability differs from the proportion of Democrats in the whole population.
\end{enumerate}

\medskip
\solution
\begin{enumerate}[(a)]
\item Before the move, $\Prob(B)=\Prob(B^c)=1/2$. Hence
\[
\Prob(D\mid B)=\frac{60/200}{1/2}=\frac35,\qquad
\Prob(D\mid B^c)=\frac{30/200}{1/2}=\frac3{10},\qquad
\Prob(D)=\frac{90}{200}=\frac9{20}.
\]
The move does not change $\Prob(D)$. Nobody changes parties, so the total
number of Democrats remains 90. The population remains 200, and selection
is uniform over these 200 residents. Thus the overall probability remains $9/20$.

\item After the move, Blue has 90 residents, including $60-k$ Democrats,
and Red has 110 residents, including $30+k$ Democrats. Therefore,
\[
\Prob(D\mid B)=\frac{60-k}{90},\qquad
\Prob(D\mid B^c)=\frac{30+k}{110}.
\]
Both proportions increase strictly if and only if
\[
\frac{60-k}{90}>\frac35\iff k<6,
\qquad
\frac{30+k}{110}>\frac3{10}\iff k>3.
\]
Thus the possible integer values are $\boxed{k=4\text{ or }5}$.
By the law of total probability, using the post-move population weights,
\[
\Prob(D)
=\Prob(D\mid B)\Prob(B)+\Prob(D\mid B^c)\Prob(B^c)
=\frac{60-k}{90}\frac{90}{200}
+\frac{30+k}{110}\frac{110}{200}
=\frac9{20}.
\]
Although both state proportions increase, the weights change from
$(1/2,1/2)$ to $(9/20,11/20)$. The higher-proportion state, Blue, receives
less weight, while Red receives more weight, leaving the overall proportion unchanged.

\newpage
\item With $k=5$, Blue contains 55 Democrats among 90 residents, and Red
contains 35 Democrats among 110 residents. The list contains all 90 Blue
residents and the 35 Democrats in Red, for a total of 125 residents.
The two conditional group weights are
\[
\Prob(B\mid A)=\frac{90}{125}=\frac{18}{25},\qquad
\Prob(B^c\mid A)=\frac{35}{125}=\frac7{25}.
\]
Within these groups, $\Prob(D\mid A,B)=55/90=11/18$ and
$\Prob(D\mid A,B^c)=1$. By the law of total probability conditional on $A$,
\[
\begin{aligned}
\Prob(D\mid A)
&=\Prob(D\mid A,B)\Prob(B\mid A)
+\Prob(D\mid A,B^c)\Prob(B^c\mid A)\\
&=\frac{11}{18}\frac{18}{25}+1\cdot\frac7{25}
=\boxed{\frac{18}{25}}.
\end{aligned}
\]
This exceeds the overall proportion $9/20$ because the list includes every
Democrat but excludes all 75 Republicans in Red. Conditioning on being
listed therefore increases the proportion of Democrats.
\end{enumerate}
\end{homeworkProblem}

\newpage
\begin{homeworkProblem}[2]
% Source: Fall 2025 HW3, Problem 3, PDF page 4.
A sequence of $n\ge1$ mutually independent trials is performed. Each trial results in either success or failure, and the probability of success on trial $i$ is $p_i$, where $0\le p_i\le1$. Let $A_n$ be the event that the number of successes among the first $n$ trials is even.
\begin{enumerate}[(a)]
\item Find $\Prob(A_2)$ by listing the possible outcomes, and show that
\[
\Prob(A_2)=\frac{1+(1-2p_1)(1-2p_2)}{2}.
\]
\item By considering whether the last trial succeeds or fails, obtain a recursion for $\Prob(A_n)$ for $n\ge2$. Use the recursion to prove that
\[
\Prob(A_n)=\frac{1+\prod_{i=1}^{n}(1-2p_i)}{2}.
\]
You may use mathematical induction or repeated substitution in the recursion. State the initial value explicitly.

%\emph{Hint:} After obtaining the recursion, consider $2\Prob(A_n)-1$.
\item Under the independence assumption, is $\Prob(A_n)=1/2$ if and only if at least one $p_i=1/2$? Justify your answer. Then remove the independence assumption and specify probabilities for the four possible outcomes of two trials so that $p_1=p_2=1/2$ but $\Prob(A_2)\ne1/2$.
\end{enumerate}

\medskip
\solution
\begin{enumerate}[(a)]
\item Let $S_i$ be the event that trial $i$ succeeds, and write $S$ and $F$
for success and failure in an ordered outcome. The outcomes in $A_2$ are
$SS$ and $FF$, with probabilities $p_1p_2$ and $(1-p_1)(1-p_2)$,
respectively. Thus, for all $0\le p_1,p_2\le1$,
\[
\Prob(A_2)=\Prob(SS)+\Prob(FF)=p_1p_2+(1-p_1)(1-p_2).
\]
When $0<p_1<1$, this can also be written by conditioning on the first trial:
\[
\begin{aligned}
\Prob(A_2)
&=\Prob(A_2\mid S_1)\Prob(S_1)
+\Prob(A_2\mid S_1^c)\Prob(S_1^c)\\
&=p_2p_1+(1-p_2)(1-p_1)\\
&=1-p_1-p_2+2p_1p_2
=\boxed{\frac{1+(1-2p_1)(1-2p_2)}{2}}.
\end{aligned}
\]
At $p_1=0$ or $1$, the same result follows directly from the outcome sum above.

\item If trial $n$ fails, the total number of successes is even exactly
when $A_{n-1}$ occurs. If trial $n$ succeeds, it is even exactly when
$A_{n-1}^c$ occurs. Independence and the law of total probability give
\[
\boxed{\Prob(A_n)=(1-p_n)\Prob(A_{n-1})
+p_n\bigl(1-\Prob(A_{n-1})\bigr)},\qquad n\ge2.
\]
The initial value is $\Prob(A_1)=1-p_1$, since zero successes is even.
This agrees with the claimed formula:
\[
\Prob(A_1)=1-p_1=\frac{1+(1-2p_1)}{2}.
\]
For the induction step, assume that
\[
\Prob(A_k)=\frac{1+\prod_{i=1}^{k}(1-2p_i)}{2}.
\]
Then the recursion yields
\[
\begin{aligned}
\Prob(A_{k+1})
&=(1-p_{k+1})\Prob(A_k)
+p_{k+1}\bigl(1-\Prob(A_k)\bigr)\\
&=p_{k+1}+(1-2p_{k+1})\Prob(A_k)\\
&=p_{k+1}+(1-2p_{k+1})
\frac{1+\prod_{i=1}^{k}(1-2p_i)}{2}\\
&=\frac{1+\prod_{i=1}^{k+1}(1-2p_i)}{2}.
\end{aligned}
\]
Therefore, by mathematical induction, the formula holds for every $n\ge1$.

\item Under independence, the formula in part (b) gives
\[
\Prob(A_n)=\frac12
\iff \prod_{i=1}^{n}(1-2p_i)=0
\iff p_i=\frac12\text{ for at least one }i.
\]
The last equivalence holds because a finite product is zero if and only
if at least one factor is zero.

Without independence, consider the joint distribution
\[
\begin{array}{c|cccc}
\text{Outcome}&SS&SF&FS&FF\\\hline
\text{Probability}&1/2&0&0&1/2
\end{array}
\]
Then $p_1=p_2=1/2$, but
\[
\Prob(A_2)=\Prob(SS)+\Prob(FF)=1\ne\frac12.
\]
The trials are dependent: $\Prob(SS)=1/2\ne p_1p_2=1/4$.
\end{enumerate}
\end{homeworkProblem}

\newpage
\begin{homeworkProblem}[3]
% Source: Fall 2025 HW2, Problem 1(a)-(b), PDF page 2.
To battle against spam, Bob installs two anti-spam programs. An email is either legitimate (event $L$) or spam (event $L^c$). Let $M_j$ be the event that program $j$ marks the email as legitimate, for $j=1,2$. Assume that
\[
\Prob(L)=0.10,\qquad
\Prob(M_j\mid L)=\Prob(M_j^c\mid L^c)=0.90
\quad(j=1,2).
\]
The two programs process the same email. Their outputs are conditionally independent given $L$, and also conditionally independent given $L^c$.
\begin{enumerate}[(a)]
\item Find $\Prob(L\mid M_1)$ and simplify your answer.
\item Find $\Prob(L\mid M_1,M_2)$ and $\Prob(L\mid M_1,M_2^c)$. Compare each with the prior probability $\Prob(L)$.
\item Are $M_1$ and $M_2$ independent without conditioning on the email's true status? Prove your answer by computing the required probabilities. Explain why your answer is consistent with the stated conditional independence.
\end{enumerate}

\medskip
\solution
\begin{enumerate}[(a)]
\item First, taking complements gives
\[
\Prob(M_j\mid L^c)=1-\Prob(M_j^c\mid L^c)=0.10,
\qquad j=1,2.
\]
By the law of total probability,
\[
\Prob(M_1)=\Prob(M_1\mid L)\Prob(L)
+\Prob(M_1\mid L^c)\Prob(L^c)
=0.90\cdot0.10+0.10\cdot0.90=0.18.
\]
Thus Bayes' rule gives
\[
\Prob(L\mid M_1)
=\frac{\Prob(M_1\mid L)\Prob(L)}{\Prob(M_1)}
=\frac{0.90\cdot0.10}{0.18}
=\boxed{\frac12}.
\]

\item Using conditional independence and the law of total probability,
\[
\Prob(M_1\cap M_2)
=0.90^2\cdot0.10+0.10^2\cdot0.90=0.09.
\]
Therefore, by Bayes' rule,
\[
\Prob(L\mid M_1,M_2)
=\frac{\Prob(M_1\cap M_2\mid L)\Prob(L)}{\Prob(M_1\cap M_2)}
=\frac{0.90^2\cdot0.10}{0.09}
=\boxed{\frac9{10}}.
\]
Similarly, conditional independence also applies when the second event
is replaced by its complement, so
\[
\Prob(M_1\cap M_2^c)
=0.90\cdot0.10\cdot0.10+0.10\cdot0.90\cdot0.90=0.09.
\]
Hence
\[
\Prob(L\mid M_1,M_2^c)
=\frac{\Prob(M_1\cap M_2^c\mid L)\Prob(L)}{\Prob(M_1\cap M_2^c)}
=\frac{0.90\cdot0.10\cdot0.10}{0.09}
=\boxed{\frac1{10}}.
\]
When both programs mark the email as legitimate, the posterior $0.90$
exceeds the prior $\Prob(L)=0.10$. When the programs disagree, the posterior
equals the prior.

\item By the law of total probability, for $j=1,2$,
\[
\Prob(M_j)=0.90\cdot0.10+0.10\cdot0.90=0.18.
\]
However, part (b) gives
\[
\Prob(M_1\cap M_2)=0.09
\ne 0.18^2=0.0324=\Prob(M_1)\Prob(M_2).
\]
Thus $M_1$ and $M_2$ are not independent without conditioning.
This is consistent with conditional independence: within each fixed email
status, the outputs are independent, but both programs are more likely
to mark a legitimate email as legitimate than a spam email. Averaging
over the shared, unknown status introduces dependence between their outputs.
\end{enumerate}
\end{homeworkProblem}

\newpage
\begin{homeworkProblem}[4]
% Mathematical source: Fall 2025 HW2, Problem 4(b)-(c), PDF page 5, with n=2.
% Extensions: optimal selective publication and the role of independence.
Consider a simplified model of an AI evaluation. In one run, an AI system answers two questions. Write $C$ for correct and $I$ for incorrect; the ordered outcomes $CC,CI,IC,II$ also denote the corresponding events. In parts (a) and (b), each answer is correct with probability $1/2$, independently of the other, so the four outcomes are equally likely.

An evaluator checks both answers against an answer key and may publish one of the following reports. Let $M_C$ and $M_I$ denote the events that the corresponding reports are published:
\[
\begin{array}{cl}
M_C:&\text{``At least one answer is correct'';}\\
M_I:&\text{``At least one answer is incorrect''.}
\end{array}
\]
At most one report is published per run, and every published report must be true.
\begin{enumerate}[(a)]
\item Initially, every run is reported. For $CC$, the evaluator publishes $M_C$; for $II$, the evaluator publishes $M_I$. For either $CI$ or $IC$, a fair coin is used to choose between the two reports, each with probability $1/2$. Using Bayes' rule, find
\[
\Prob(CC\mid M_C)\qquad\text{and}\qquad\Prob(II\mid M_I).
\]
\item Replace the rule in part (a) by any randomized reporting rule based on the outcome. The evaluator may choose not to publish a report for any of the four outcomes. The two reports must each have positive probability and must satisfy
\[
\Prob(CC\mid M_C)\ge\frac23,\qquad
\Prob(II\mid M_I)\ge\frac34.
\]
What is the smallest possible probability that no report is published? Specify a reporting rule for each of $CC,CI,IC,II$ that attains your answer, and prove that no admissible rule can do better.
\item Remove the independence assumption, but keep each answer's probability of being correct equal to $1/2$. The evaluator may choose a new randomized reporting rule, but must publish exactly one truthful report per run, with both report types having positive probability. Let $t=\Prob(CC)$. Determine all values of $t$ for which such a rule satisfies both posterior inequalities in part (b). Prove necessity and sufficiency, giving a valid rule for every claimed value of $t$. The rule may depend on $t$.
\end{enumerate}

\medskip
\solution
\begin{enumerate}[(a)]
\item By the law of total probability,
\[
\Prob(M_C)=1\cdot\frac14+\frac12\cdot\frac14
+\frac12\cdot\frac14=\frac12.
\]
Similarly, $\Prob(M_I)=1/2$. Bayes' rule therefore gives
\[
\Prob(CC\mid M_C)=\frac{\Prob(M_C\mid CC)\Prob(CC)}{\Prob(M_C)}
=\frac{1\cdot(1/4)}{1/2}=\boxed{\frac12},
\]
\[
\Prob(II\mid M_I)=\frac{\Prob(M_I\mid II)\Prob(II)}{\Prob(M_I)}
=\frac{1\cdot(1/4)}{1/2}=\boxed{\frac12}.
\]

\newpage
\item Let $N$ be the event that no report is published. The posterior constraints imply
\[
\Prob(M_C)\le\frac32\Prob(CC\cap M_C)
\le\frac32\Prob(CC)=\frac38,
\]
\[
\Prob(M_I)\le\frac43\Prob(II\cap M_I)
\le\frac43\Prob(II)=\frac13.
\]
Since the reports are mutually exclusive,
\[
\Prob(N)=1-\Prob(M_C)-\Prob(M_I)
\ge1-\frac38-\frac13=\frac7{24}.
\]
To attain this bound, use the following rule. Each row gives the conditional
probabilities of the three reporting choices given the outcome.
\[
\begin{array}{c|ccc}
\text{Outcome}&M_C&M_I&\text{No report}\\\hline
CC&1&0&0\\
CI&1/4&1/6&7/12\\
IC&1/4&1/6&7/12\\
II&0&1&0
\end{array}
\]
Every published report is true. For this rule,
\[
\Prob(M_C)=\frac14+\frac12\cdot\frac14=\frac38,
\qquad
\Prob(M_I)=\frac14+\frac12\cdot\frac16=\frac13.
\]
Thus both report types have positive probability, and
\[
\Prob(CC\mid M_C)=\frac{1/4}{3/8}=\frac23,
\qquad
\Prob(II\mid M_I)=\frac{1/4}{1/3}=\frac34.
\]
Finally, $\Prob(N)=(1/2)(7/12)=7/24$. Hence the smallest possible
probability of publishing no report is $\boxed{7/24}$.

\newpage
\item The marginal probabilities of correctness imply
\[
\Prob(CI)=\frac12-t,\qquad
\Prob(IC)=\frac12-t,\qquad
\Prob(II)=t.
\]
Nonnegativity requires $0\le t\le1/2$.

\textbf{Necessity.}
Exactly one truthful report must be published per run. Consequently,
$CC$ must produce $M_C$, and $II$ must produce $M_I$, so
\[
\Prob(CC\cap M_C)=\Prob(II\cap M_I)=t.
\]
The posterior constraints give
\[
\Prob(M_C)\le\frac32t,\qquad
\Prob(M_I)\le\frac43t.
\]
Since $\Prob(M_C)+\Prob(M_I)=1$, it follows that
\[
1\le\left(\frac32+\frac43\right)t=\frac{17}{6}t,
\qquad\text{so}\qquad t\ge\frac6{17}.
\]

\textbf{Sufficiency.}
For any $6/17\le t\le1/2$, use the rule
\[
\begin{array}{c|cc}
\text{Outcome}&M_C&M_I\\\hline
CC&1&0\\
CI&3/5&2/5\\
IC&3/5&2/5\\
II&0&1
\end{array}
\]
Each row gives conditional reporting probabilities. The rule always
publishes exactly one truthful report, and
\[
\Prob(M_C)=t+\frac35(1-2t)=\frac{3-t}{5},\qquad
\Prob(M_I)=t+\frac25(1-2t)=\frac{2+t}{5}.
\]
Both probabilities are positive throughout the claimed interval. Moreover,
\[
\Prob(CC\mid M_C)=\frac{5t}{3-t}\ge\frac23
\iff 17t\ge6,
\]
\[
\Prob(II\mid M_I)=\frac{5t}{2+t}\ge\frac34
\iff 17t\ge6.
\]
Thus both requirements hold, including at the endpoints. The set of all
feasible values is therefore $\boxed{t\in[6/17,\,1/2]}$.
\end{enumerate}
\end{homeworkProblem}

\end{document}
